\section{Discussion of Findings}
\label{sec:results_discussion_of_findings}

The experiments show that GLPE is most effective when the agent must discover useful behavior before extrinsic rewards provide sufficient guidance.
This conclusion is supported by \textsc{MountainCar-v0}, where GLPE achieves the best final return, strongest step-AUC, and complete solved-threshold reliability.
The result is consistent with the design of the intrinsic reward: feature-space impact encourages behaviorally meaningful transitions, while region-local learning progress favors transitions whose dynamics are becoming more predictable.

The results also show that gating is not universally beneficial.
The gate improves the sparse-reward result but can reduce performance in dense-reward or high-variance environments.
This is most visible in \textsc{Humanoid-v5}, where gated GLPE has lower final performance and lower threshold reliability than the ungated variant.
A plausible explanation is that high-dimensional locomotion produces noisy local progress estimates, causing the gate to suppress regions that are difficult but still useful for learning.
Thus, the empirical role of the gate is best described as a selective guardrail against unproductive curiosity, not as a general-purpose performance booster.

The ablation results indicate that the two reward components contribute differently across tasks.
Impact-only exploration is competitive in \textsc{Ant-v5}, learning-progress-only exploration is strong in \textsc{BipedalWalker-v3}, and the full combination is necessary for the strongest \textsc{MountainCar-v0} result.
This pattern suggests that the fixed component weights used in the experiments are a reasonable general setting but are unlikely to be optimal for every environment.
An adaptive weighting mechanism could potentially improve robustness by increasing the importance of learning progress in sparse-reward settings and increasing the importance of impact in dense-control settings.

The threshold analyses reveal that partial competence and solved performance should be distinguished.
GLPE and GLPE without gating reach 25\% and 50\% thresholds reliably in most environments, but full solved-threshold reliability is more difficult in \textsc{BipedalWalker-v3} and \textsc{Humanoid-v5}.
This distinction is important because mean learning curves can make methods appear similar even when their probability of reaching a practical target differs across seeds.
Reliability metrics therefore provide information that is not captured by mean return alone.

Wall-clock results show that computational overhead must be considered when evaluating intrinsic-reward methods.
The gated method adds region statistics and robust threshold computations, and these costs matter most when environment simulation is inexpensive.
The ungated variant often provides a better runtime-performance balance because it retains the combined impact--learning-progress reward without the additional gate computation.
For practical use, the choice between gated and ungated GLPE should therefore depend on whether the expected benefit of suppressing unproductive regions is likely to exceed the added runtime cost.

Overall, the findings support the central premise of this study: learning progress can be combined with feature-space impact to produce an intrinsic reward that is useful for exploration, particularly in sparse-reward settings.
At the same time, the results identify limitations in high-dimensional and dense-reward environments, where the gate may become conservative and where standard baselines remain competitive.
The method should therefore be viewed as a contribution to targeted intrinsic motivation rather than as a universally dominant reinforcement-learning algorithm.